\section{A class-number shortcut to finite-field descent fails}
\label{root-class-number-descent}

\textbf{Scope.} This source check concerns one exact assertion in
Thakur's January 2020 preprint, not the other results of that paper.
It uses the classical abelian-variety corollary of CC-AG14; no
additional theorem about higher-codimension abelian cycles, and no
adopted analytic Riemann hypothesis, enters the counterexample.

\subsection{BR-15-CLASSBOUND --- Refutation of the stated class-number bound}

Thakur (2020, Definition 2.3 and Proposition 2.9, p.~6,
arXiv:2001.00210v1) defines a model by isogeny after base extension
and asserts that a type $\mathrm{IV}(1,d)$ abelian variety has such a
model over $\mathbf F_{p^{h_K}}$, where $h_K$ is the class number of
the centre $K$ of its geometric endomorphism algebra. The assertion
is false, including if one weakens the demanded isogeny to a geometric
isogeny.

\textbf{Counterexample and complete deduction.} Let
\[
 K=\mathbf Q(s),\quad s^2=-23,\quad
 \omega=(1+s)/2,\quad \eta=1+\omega,\quad
 \gamma=8\bar\eta=12-4s.
\]
By the fully proved classical corollary of CC-AG14, there is a
geometrically simple abelian threefold $A/\mathbf F_{2^9}$ with
Frobenius polynomial
\[
                 (T^2-24T+512)^3,
\]
and no abelian variety over $\mathbf F_8$ is geometrically isogenous
to $A$. The corollary proves this using the classical Honda--Tate
classification and its local-index formula: every possible source
label giving a single geometric copy has degree three over $K$ and
would need index one, whereas the complete root-of-unity and local
cube calculation forces index three. Thus its nonexistence conclusion
allows every source label and every root-of-unity twist; it is not
only a failure to take a particular cube root of $\gamma$.

We verify the remaining class-number hypothesis explicitly. The ring
of integers is $\mathcal O_K=\mathbf Z[\omega]$, its discriminant is
$-23$, and
\[
 N_{K/\mathbf Q}(x+y\omega)=x^2+xy+6y^2
             =(x+y/2)^2+23y^2/4.
\]
The maximality of this quadratic order follows because its
discriminant is squarefree: the square of the index of any larger
order would divide $23$. Put
\[
 \mathfrak p=(2,\omega-1),\quad
 \bar{\mathfrak p}=(2,\omega),\quad
 \mathfrak q=(3,\omega),\quad
 \bar{\mathfrak q}=(3,\omega-1).
\]
The polynomial of $\omega$ is $T^2-T+6$; its distinct roots modulo
$2$ and modulo $3$ give precisely these primes, each with residue
degree one. In particular
$(2)=\mathfrak p\bar{\mathfrak p}$ and
$(3)=\mathfrak q\bar{\mathfrak q}$.
The norm of $\eta$ is $8$, and its residues at the two primes above
$2$ show $(\eta)=\mathfrak p^3$. No element has norm $2$: if $y\ne0$
the displayed quadratic form is at least $23/4$, and if $y=0$ it is
a square. Hence $\mathfrak p$ is not principal and its class has
order exactly three. Also $N(\omega)=6$, and its residues give
\[
                  (\omega)=\bar{\mathfrak p}\mathfrak q,
             \qquad [\mathfrak q]=[\mathfrak p].
\]
Minkowski's ideal-class bound (Milne, 2020, Theorem 4.3, p.~70)
gives an integral representative of norm at most
$2\sqrt{23}/\pi<4$ for every class. An ideal of norm $2$ or $3$ is
prime, so the above four primes, together with the unit ideal,
exhaust these representatives. Their classes are
$1,[\mathfrak p],[\mathfrak p]^{-1}$. Therefore $h_K=3$.

The endomorphism division algebra of $A$ has centre $K$, degree
three, and is ramified only at the two primes above $2$, with
invariants of denominators three. These statements are part of the
Honda--Tate computation in CC-AG14. They also describe the geometric
endomorphism algebra: over every finite extension the centre and
index remain unchanged, and the injective maps between the
endomorphism algebras therefore have the same finite dimension and
are isomorphisms. Every geometric endomorphism descends to a finite
extension. A polarization gives the required involution of the
second kind. Thus $A$ satisfies Thakur's Definition 2.1, p.~3, for
type $\mathrm{IV}(1,3)$.

The asserted bound would now give a model over
$\mathbf F_{2^{h_K}}=\mathbf F_8$, which is a subfield of the
displayed field of definition $\mathbf F_{2^9}$. Any isogeny over
$\mathbf F_{2^9}$ would give a geometric isogeny, contradicting
CC-AG14. This proves the refutation in the source's own model
convention. $\square$

\subsection{Location of the failed inference}

The proof of the cited Proposition 2.9 chooses a generator $\pi_0$
of $\mathfrak p^{h_K}$ and assigns its associated abelian variety
dimension $d$ and the original division algebra. When $p$ splits
in the imaginary quadratic field, this generator instead has slopes
$1$ and $0$. The classical local-index formula gives index one;
its centre is $K$, since its two valuations are different. Its
Honda--Tate variety is therefore an ordinary elliptic curve.
For $d>1$, it cannot supply the asserted variety. This is the
precise failed inference; the argument above independently
refutes the conclusion rather than relying only on a gap in that
proof.

\subsection{BR-15-SLOPEBOUND --- The corrected geometric construction}

For clarity, the slope-sensitive construction in Thakur (2020,
Proposition 2.4, pp.~4--5) has a different exponent. Here is the
exact form needed for geometric descent.

Let $A$ be a geometrically simple abelian variety over a finite field
of characteristic $p$ whose geometric endomorphism division algebra has imaginary
quadratic centre $K$ and degree $d>1$. Suppose $p$ splits in $K$,
and write its invariant at one prime $\mathfrak p$ as $j/d$ with
$0<j<d$ and $\gcd(j,d)=1$. Let $t$ be the order of
$[\mathfrak p]$ in $\operatorname{Cl}(K)$. Then the geometric
isogeny class of $A$ has a representative over
$\mathbf F_{p^{td}}$, hence also over $\mathbf F_{p^{h_Kd}}$.
This is a sufficient geometric field, not a claimed smallest field
and not a prescribed finite-field Frobenius descent.

\textbf{Complete proof of this form.} Choose a finite field
$\mathbf F_{p^N}$ defining $A$ and a finite $\mathbf Q$-basis of
its geometric endomorphism algebra. Such a field exists because each
basis element, after clearing a denominator, is a morphism defined
over some finite extension. Thus the rational endomorphism algebra
over this field is already the full geometric algebra.
Frobenius $\pi$ commutes with those endomorphisms, so $\pi\in K$;
the Honda--Tate centre formula gives $\mathbf Q(\pi)=K$.
At the split primes the local degree is one. Its slopes lie in
$[0,1]$ and represent the nonzero invariants $j/d$ and
$(d-j)/d$, so they equal those fractions themselves. In particular
$d\mid N$ and $\dim A=d$.

Choose $\eta\in\mathcal O_K$ with $(\eta)=\mathfrak p^t$.
Its norm is $p^t$. Put
\[
                  \rho=\eta^j\bar\eta^{\,d-j}.
\]
It is an integral Weil $p^{td}$ number with the same two slopes.
The fractions cannot both be $1/2$: if they were, the two valuations
of $\pi$ would be equal, and $\pi^2/p^N$ would be a unit of
the imaginary quadratic field. Its finite order would make a power
of $\pi$ rational, contradicting the unchanged geometric centre
$K$ over every finite extension where all endomorphisms are already
defined. Hence $j\ne d-j$, and $\rho$ and all its positive powers
have unequal valuations and centre $K$. Their index is $d$.
Honda--Tate therefore gives an abelian variety $A_0$ of dimension
$d$ over $\mathbf F_{p^{td}}$ with Frobenius $\rho$.

The elements $\rho^N$ and $\pi^{td}$ have identical valuations
at every prime of $K$, so their ratio is a unit. The unit group of
an imaginary quadratic field is finite (equivalently, its lattice
has only finitely many points on the unit circle). If this ratio
has order $e$, then
\[
                         \rho^{Ne}=\pi^{tde}.
\]
Over $\mathbf F_{p^{Ntde}}$ the two Frobenius eigenvalues and their
conjugates thus agree, with the same multiplicity $d$. Tate's
isogeny criterion makes the two base extensions isogenous. This
gives the geometric representative $A_0$. Finally $t\mid h_K$,
so extension of $A_0$ to $\mathbf F_{p^{h_Kd}}$ is available.
$\square$

The exact finite divisor profile of CC-AG17 remains stronger than
this sufficient construction. The present source correction is not
counted as a reduction of the remaining arbitrary-variety cycle gate.

\subsection{Primary locators}

Thakur, S. (2020). \emph{Isogenies of certain abelian varieties over
finite fields with $p$-ranks zero} (arXiv:2001.00210v1, January 1).
https://arxiv.org/abs/2001.00210v1 . Definition 2.1, p.~3;
Proposition 2.4, pp.~4--5; Definition 2.3 and Proposition 2.9 with
its proof, p.~6. The refutation refers exactly to this version.

Milne, J. S. (2020). \emph{Algebraic number theory} (Version 3.08,
July 19). Theorem 4.3, p.~70, ideal-class Minkowski bound.
https://www.jmilne.org/math/CourseNotes/ANT.pdf .

Tate, J. (1969). Classes d'isog\'enie des vari\'et\'es ab\'eliennes
sur un corps fini (d'apr\`es T.~Honda). \emph{S\'eminaire Bourbaki},
Expos\'e 352, pp.~95--110. Theorem 1(i)--(ii) and Remark 2,
p.~96. https://www.numdam.org/article/SB_1968-1969__11__95_0.pdf .

Tate, J. (1966). Endomorphisms of abelian varieties over finite
fields. \emph{Inventiones Mathematicae, 2}, 134--144.
Theorem 1(c), p.~139, Frobenius-polynomial isogeny criterion.
https://doi.org/10.1007/BF01404549 .
